% sections/01_introduction.tex: the question, why it is hard, the approach, the status note, related work.
% Numbers only through \R / \Rci. Static table: tab:related (prior work; literature only, no result of ours).
% The status note is the ONLY place that states the audit status of the final-model lane.
\section{Introduction}\label{sec:intro}

\paragraph{The question.} A person holds a smartphone at arm's length and brings it to touch the nose, with the left
or the right hand. The task, set by the principal investigator (PI), is to predict which hand from the phone's
motion sensors alone, as the browser's DeviceMotion interface delivers them to a web page at about
60~Hz~\cite{w3c2025deviceorientation}. The answer is a percentage, \pright{}: the probability that the right hand
moved the phone. It should be calibrated (of the touches given 80\,\%, about 80\,\% should be right-hand touches),
and it must hold for persons the model has never seen, because the PI will test it on new persons.

\paragraph{Why it is hard.} Five things make the task hard.
\begin{enumerate}[nosep,leftmargin=1.5em]
  \item Persons differ in arm length, speed and style, and the model sees no movement of the test person.
  \item Phones differ in sensors, sampling and browser~\cite{stisen2015heterogeneity}: iOS reports the
    accelerometer with the opposite sign to Android, and the browser delivers only about 60~Hz.
  \item The device axes turn with the grip, so the same arm movement gives different device-frame signals when the
    phone is held upright, sideways or with the screen away.
  \item The two hands make mirror-image movements, so most simple descriptors (speed, duration, path length) are
    the same for both hands.
  \item In production there is no start cue: the page must find the approach in a continuous stream, reject other
    movements such as the return of the arm, and decide when it has seen enough.
\end{enumerate}

\paragraph{The approach.} The analysis rests on the left--right mirror symmetry of the body
(Section~\ref{sec:physics}). Mirroring a movement flips the sign of a few chiral (mirror-odd) quantities, chiefly
the rotation about gravity, and leaves the rest unchanged, so only chiral quantities can carry the hand. Every trial
is standardised into a gravity frame that does not depend on the grip and a body frame that follows the phone's
heading (Section~\ref{sec:methods}). Four model families, from a one-line rule to deep sequence models, are
evaluated by leave-one-person-out (LOSO) cross-validation, with all tuning nested inside the training persons. A
pre-specified cross-lane rule selects the final model among the lanes' candidates, and a nested run of this
selection gives its selection-aware estimate (Section~\ref{sec:results}). Calibration, several attempts, natural free-choice use and
robustness to an unseen grip are tested on the same folds. A separate simulation removes the start cue and adds
detection, online logic and a stopping rule (Section~\ref{sec:production}). Secondary analyses ask what else the
same data reveal: speed, grip, handedness and the value of demographics (Section~\ref{sec:secondary}). The
appendices state every decision with its options, evidence and rationale.

\paragraph{Status of this version.} This is a preliminary version, written for the researcher's review. The
analysis was organised in lanes, and the appendices keep their names. The lanes on interpretable models (m1,
Appendix~\ref{app:m1}), classical machine learning (m2, Appendix~\ref{app:m2}), time-frequency and deep models
(m3, Appendix~\ref{app:m3}), the cue-free production pipeline (m4, Appendix~\ref{app:prod}), preprocessing (p1,
Appendix~\ref{app:prep}) and the case study (x1, Appendix~\ref{app:pd01}) each passed two independent audits, one
for rigour and one for reproducibility, and were fixed after each. The final-model lane (selection, calibration and
export; Appendix~\ref{app:final}) is a first pass that has not yet been independently audited. The secondary
analyses of speed, grip and handedness passed two audits with fixes; the demographics test is a quick, unaudited
first pass. The final model has not yet been run inside the cue-free production simulation, so the production
numbers of Section~\ref{sec:production} use that simulation's own reference classifier. No review of the
report's completeness has run yet.

\paragraph{Related work.} Table~\ref{tab:related} lists prior work on detecting the hand or the handedness from
phone or wrist sensors. Most studies infer the hand from touch traces or grip during ordinary phone use
\cite{goel2012gripsense,guo2016operating,nelavelli2018adaptive,fernandez2021implicit}, and several split the data
by trial or at random, so that the same person can appear in training and test. Such splits overstate the accuracy
for new persons \cite{saeb2017usecase,little2017crossvalidation,chaibubneto2019subject}, and with few persons a
cross-validated estimate has wide error bars \cite{varoquaux2018cvfailure}. Closest to this task,
Mohammadi-Ghazi et al.\ saw opposite rotation senses of the two hands in a clinical finger-to-nose test with wrist
sensors, but built no classifier~\cite{mohammadighazi2022ataxia}. The studies in Table~\ref{tab:related} report
accuracies; a calibrated probability is a separate requirement, and modern classifiers, deep networks in
particular, are often over-confident~\cite{guo2017calibration}. This study therefore evaluates person-wise
throughout and checks calibration for every model it ships.

\begin{table}[htbp]
\centering
\caption{Prior work on detecting the hand or the handedness from phone or wrist sensors. Accuracies are as
reported; most evaluations are not person-wise.}\label{tab:related}
\footnotesize
\begin{tabularx}{\linewidth}{@{}l>{\raggedright\arraybackslash}X r >{\raggedright\arraybackslash}X >{\raggedright\arraybackslash}X@{}}
\toprule
Study & Signal & N & Result & Caveat \\
\midrule
GripSense \cite{goel2012gripsense} & gyroscope, touch, vibration & 10 & 84.26\,\% hand posture & controlled study \\
Löchtefeld et al.\ \cite{lochtefeld2015handedness} & sensors during the unlock gesture, $k$-NN & 12 & 98.51\,\% & all right-handed \\
Guo et al.\ \cite{guo2016operating} & touch traces; acc + gyro for hand changes & 14 & $\approx$94\,\% precision & random 10-fold CV \\
Nelavelli \& Ploetz \cite{nelavelli2018adaptive} & sign of the swipe-arc curvature & 12 & 99.53\,\% consistent & no held-out users \\
Fernández et al.\ \cite{fernandez2021implicit} & touch features & 174 & up to 99.92\,\% TPR & oversampling before random splits \\
Chen et al.\ \cite{chen2023handedness} & motion sensors, LSTM & n/r & 89.2--97.4\,\% & unclear if test users unseen \\
Singh \& Singh \cite{singh2025dominant} & acc, gyro, orientation; XGBoost & n/r & 85.04\,\% (3 classes) & protocol not reported \\
Mohammadi-Ghazi et al.\ \cite{mohammadighazi2022ataxia} & wrist IMU, finger-to-nose & 17 & opposite rotation sense & clinical, no classifier \\
\bottomrule
\end{tabularx}
\end{table}
